% Generated draft: requires mathematical and source review.
\documentclass[11pt]{article}
\usepackage[margin=25mm]{geometry}
\usepackage{fontspec}
\setmainfont{Latin Modern Roman}
\usepackage{amsmath,amssymb,mathtools,unicode-math}
\setmathfont{Latin Modern Math}
\usepackage{enumitem,needspace,xurl,hyperref}
\setlength{\parindent}{0pt}
\setlength{\parskip}{5pt}
\setlength{\emergencystretch}{3em}
\newcommand{\N}{\mathbb{N}}
\newcommand{\Z}{\mathbb{Z}}
\newcommand{\Q}{\mathbb{Q}}
\newcommand{\R}{\mathbb{R}}
\newcommand{\C}{\mathbb{C}}
\newcommand{\F}{\mathbb{F}}
\newcommand{\A}{\mathbb{A}}
\newcommand{\PP}{\mathbb{P}}
\newcommand{\E}{\mathbb{E}}
\newcommand{\eps}{\varepsilon}
\newcommand{\dd}{\,\mathrm{d}}
\begin{document}
\section*{UnsolvedMath ID: 30000606}
\textbf{Source problem code:} catalogue entry
\subsection{Definitions and mathematical statement}
Let \(H=(V,E)\) be a finite simple \(k\)-uniform hypergraph, with \(|V|=n\) and \(E\subseteq\binom Vk\). For \(S\in\binom V{k-1}\), define its codegree and the minimum codegree by
\[
d_H(S)=|\{v\in V\setminus S:S\cup\{v\}\in E\}|,
\qquad\delta_{k-1}(H)=\min_{S\in\binom V{k-1}}d_H(S).
\]
The source's proposed sufficient threshold is exactly
\[
\delta_{k-1}(H)\geq\left\lfloor\frac{n-k+3}{2}\right\rfloor.
\]
It asks whether this forces a Hamiltonian cycle.

A natural intended interpretation in a codegree problem is a \emph{tight} Hamiltonian cycle: a cyclic ordering \(v_1,\ldots,v_n\) of every vertex such that
\[
\{v_i,v_{i+1},\ldots,v_{i+k-1}\}\in E\qquad(i\pmod n).
\]
With a nondegenerate cycle range such as \(n\geq k+1\), this describes the desired spanning cyclic sequence of consecutive \(k\)-sets. The record does not actually specify tight rather than loose or Berge cycles; these notions are different and may have different divisibility restrictions. This display is therefore an explicit candidate meaning, not a silently established definition from the source.

The ranges of \(k\) and \(n\), any sufficiently-large-\(n\) qualifier, and small or degenerate cycle conventions are absent. Once these are fixed, the intended assertion is that every admissible \(H\) satisfying the displayed non-strict floor bound has the corresponding cycle. No rounding change, asymptotic error term, degree substitution, or missing divisibility condition is supplied by inference.

The supplied literature triage checked 2026-08-21 reports no authoritative complete resolution of the extracted claim. Because the cycle type and size regime are missing, that aggregate open label is not independently certified for each possible precise formulation.
\subsection{Short English statement}
Does the exact minimum-codegree threshold \(\lfloor(n-k+3)/2\rfloor\) force a spanning hypergraph cycle? The source must fix the cycle notion and the admissible size range; the usual tight-cycle reading is made explicit as an interpretation.
\subsection{Sources}
\begin{itemize}
\item Local UnsolvedMath registry, permanent ID 30000606, catalogue entry, \emph{Codegree Thresholds for Hypergraph Hamilton Cycles}.
\item \emph{Combinatorics, Probability and Computing} (2007), Oberwolfach Reports, pp. 2875--2952, DOI \url{https://doi.org/10.4171/owr/2006/48}.
\item Supplied source-verification claim and literature triage checked 2026-08-21 cite only the original report; no exact cycle definition, size hypothesis, or separate threshold theorem is supplied.
\end{itemize}
\end{document}
